\clearpage
\fancyhead[L]{\small GAUSS AN VON LINDENAU.}
\hypertarget{native-316-lower}{}
\begin{center}
\textbf{GAUSS AN VON LINDENAU.}

\par\vspace{4pt}\hrule\vspace{4pt}
\fontsize{10}{12}\selectfont Monatliche Correspondenz. Band XIX. Seite 501..503. 1809. Mai.
\par\vspace{4pt}\hrule\vspace{4pt}
\end{center}

\begin{flushright}Göttingen, am 20. Mai 1809.\end{flushright}

\par\indent Ich schicke Ihnen beiliegend die Ergebnisse meiner Rechnungen zur Verbesserung der Elemente der neuen Planeten. Einen Teil dieser Rechnungen habe ich schon vor einiger Zeit abgeschlossen. Bei Ceres werden Sie einige Unterschiede zwischen meinen und \textsc{Santini}s Ergebnissen (Monatl. Corresp. Febr. S. 191) bemerken. Wahrscheinlich liegt das zum Teil daran, dass \textsc{Santini} beim Vergleich der \emph{einzelnen} Beobachtungen mit den XI. Elementen die Störungen nicht berücksichtigt hat. Das kann ich nicht billigen, denn diese Störungen haben im Lauf der 22 Tage keineswegs gleichbleibenden Einfluss auf die geozentrischen Positionen. Die mittlere Differenz darf daher nicht als die wahre gelten. Auch bei der anschließenden Berechnung der Störungen scheint Herr \textsc{Santini} einige Fehler gemacht zu haben. Nach meinen Tafeln finde ich für die Störung der Länge $+6'\;6''4$ und für die Störung des Radiusvektors $+0{,}001570$. Sie werden bemerken, dass diese Zahlen mit denen \textsc{Santini}s fast keine Ähnlichkeit haben. \SourceLeaf{317}Die Störung der Breite, die sich freilich nicht mit der von \textsc{Santini} vergleichen lässt, weil ich mit vollständigeren Gleichungen gerechnet habe, beträgt $+56''8$. — Auf dem beiliegenden Blatt schicke ich Ihnen zugleich die Ephemeriden für die nächste Sichtbarkeit von Ceres, Juno und Vesta. Die Ephemeride für Juno hat Herr Professor \textsc{Harding} nach den bereits früher in der Monatl. Corresp. abgedruckten Elementen berechnet. Die Ephemeride für Vesta nach den IV Elementen stammt von Herrn Dr. \textsc{Schumacher}, der sich jetzt bei uns aufhält. Die für Ceres habe ich selbst übernommen. [\textsc{Von Lindenau} hat diese Ephemeride in den Aufsatz „Fortgesetzte Nachrichten u.\,s.\,f.“ im Maiheft 1809 der Mon. Corr., Seite 504 bis 515, aufgenommen. Sie gibt für Mitternacht in Göttingen vom 18. Juli 1809 bis zum 16. März 1810 in Abständen von vier Tagen die geozentrische Rektaszension und Deklination auf die Minute genau und den Logarithmus des Abstands mit vier Dezimalstellen an.] Es fehlt also nur noch Pallas. Leider habe ich von ihr bisher keine brauchbaren Beobachtungen aus dem vergangenen Jahr zu Gesicht bekommen und konnte daher noch keine Rechnungen für sie anstellen.

\Needspace{23\baselineskip}
\par\indent Zur Anzeige meines Programms im Februarheft [in der ein nicht genannter Verfasser „in dieser Zeitschrift die wichtigsten Formeln, auf denen die Methode beruht, kurz dargestellt“ und die Nummerierung der Formeln und die Bezeichnungsweise aus der Schrift von \textsc{Gauss}, Methodus peculiaris elevationem poli determinandi, Seite 39 und die folgenden Seiten dieses Bandes, unverändert übernommen hat] muss ich noch bemerken, dass dort irrtümlich von gleichen Höhen die Rede ist (S. 136). Eine solche Bedingung gehört überhaupt nicht zur Aufgabe. Auch die folgende Ergänzung ist angebracht. „Bestimmt man $\lambda$ nach Formel 9 aus der Tangente, so bleibt offen, ob $\lambda$ im ersten oder zweiten Halbkreis zu wählen ist. Formel 4 beseitigt diese Mehrdeutigkeit. Sie zeigt, dass $\lambda$ in demselben Halbkreis zu wählen ist, in dem $u$ liegt, weil $\cos\varphi$ und $\cos h$ ihrer Natur nach offensichtlich positive Größen sind.“

\par\indent Natürlich war dieser Zusatz nur der analytischen Vollständigkeit halber nötig. Welcher Wert von $\lambda$ der richtige ist, kann kaum zweifelhaft sein, denn beim zweiten Wert läge der Stand der Uhr um 12 Stunden daneben, und die Polhöhe hätte das falsche Vorzeichen. Auch aus Ihrem Auszug kann sich jeder leicht die geometrischen Erörterungen ergänzen, die das Programm selbst enthält, und zeigen, was die verwendeten Zeichen auf der Sphäre bedeuten. Herr Professor \textsc{Harding} hat übrigens für das astronomische Jahrbuch eine deutsche Übersetzung davon veranlasst.
